\section{Conclusion}\label{sec:conclusion}
At normalization exactly one, the query cost of converting $H$ to $I-H$ depends on the spectrum and the required accuracy. Known two-point spectra can admit exact constant-query conversion; a low band forces a sequence of optimal query powers determined by globally nonnegative polynomial approximation. Explicit Chebyshev thresholds and pinned constructions determine the optimal order for every fixed $K$, including at the thresholds. Requiring the polynomial to be even changes this staircase: $F_1=F_2$ gives exponent $5/6$ at a concrete accuracy where unrestricted conversion has exponent $1/2$. When $K\leq\delta^\beta$ for a fixed $\beta>0$, the cost is instead $\Theta(\delta^{-1}\log(1/K))$.

The main open question is the optimal cost, up to constant factors, in the intermediate regime. Away from the upper tier boundaries, our bounds differ by at most a factor quadratic in the threshold index; near them, one must also control the margin factor $(G_{r-1}-K)^{1/r}$. Explicit formulas for the higher even thresholds $F_r$ would sharpen the parity comparison. A separate problem is the precision and classical computation needed to synthesize the threshold-attaining circuits, particularly at a threshold ($K=G_r$, or $K=F_r$ in the even staircase), where the required accuracy $K\delta$ equals their error bound.

The exactly central LP shows that the input access model and the output task must be stated explicitly. A dual Newton state, a reusable complement block, and a digital description of the matrix are different resources, even for one sparse instance family. Our separation concerns conversion from the specified block encodings; whether a larger algorithm can avoid this conversion, or obtain stronger access efficiently, needs separate analysis.
